\begin{statement}{E.39}{Proved}{Full periodic-braid conjugacy, including the order-four test}
Let $n\ge3$, and use the actual based braid action
\textup{(E.30.4)} of $F=(\lambda,f)\in\widehat{\rm GT}$.
In $\widehat B_n$ put
\[
 \delta_n=\sigma_1\cdots\sigma_{n-1},\qquad
 \varepsilon_n=\sigma_1\delta_n,\qquad z_n=\delta_n^n.
 \tag{E.39.1}
\]
There are $h_\delta,h_\varepsilon\in\widehat B_n$ such that
\[
 F(\delta_n)=h_\delta\delta_n^\lambda h_\delta^{-1},
 \qquad
 F(\varepsilon_n)=h_\varepsilon\varepsilon_n^\lambda
                                  h_\varepsilon^{-1}.
 \tag{E.39.2}
\]
The same assertion holds for every discrete periodic braid,
by its actual conjugacy to a power of one of these two elements.

In particular, for any finite quotient $q:\widehat B_4\to T$,
write $a=q(\sigma_1)$, $b=q(\sigma_2)$, $c=q(\sigma_3)$
and $a'=qF(\sigma_1)$, $b'=qF(\sigma_2)$, $c'=qF(\sigma_3)$.
Then the following conjugacies hold in $T$ itself:
\[
 a'b'c'\sim_T(abc)^\lambda,
 \qquad a'^2b'c'\sim_T(a^2bc)^\lambda.
 \tag{E.39.3}
\]
This gives a necessary full-GT test on each actual finite
shadow.  The first equation includes the order-four
periodic element in the central quotient; it also implies
the order-two test after squaring.  No invariance of
$\ker q$ under a finite shadow is assumed.
\end{statement}
\begin{proof}
We specify both completion and orientation comparisons.
Cap the boundary of the $n$-punctured disk by a disk
whose center is the additional marked point $n+1$.
Extension of homeomorphisms by the identity gives
\[
 B_n/\langle z_n\rangle
       \xrightarrow{\sim}
 \operatorname{Stab}_{\Gamma_{0,[n+1]}}(n+1).
 \tag{E.39.4}
\]
The kernel before taking the quotient is the integer
boundary-twist group, generated by the full twist $z_n$.
The stabilizer has index $n+1$.  Its intrinsic profinite
topology agrees with the topology induced from the full
mapping class group: for each finite-index subgroup of
the stabilizer, its core in the full group has finite
index and is contained in it.  Thus completion gives
the corresponding isomorphism onto the actual open
stabilizer in $\widehat\Gamma_{0,[n+1]}$.
The kernel of $\widehat B_n$ onto the completed quotient
is the closure of $\langle z_n\rangle$.  The exponent-sum
map $B_n\to\mathbb Z$ sends $z_n$ to $n(n-1)$;
its continuous extension proves that this closure is
isomorphic to $\widehat\mathbb Z$.
Consequently
\[
 \widehat B_n/\langle z_n\rangle
       \xrightarrow{\sim}
 \operatorname{Stab}_{\widehat\Gamma_{0,[n+1]}}(n+1)
 \tag{E.39.5}
\]
uses the actual closed procyclic center.

The capping map is equivariant for the specified action.
Precisely, apply the generator formulas (1.1) in
Nakamura--Schneps \cite{BCNakamuraSchneps2000}
along $\widehat B_n\to\widehat B_{n+1}$ and then kill
the full twist and Hurwitz element in $\widehat B_{n+1}$.
Their power formulas preserve both closed normal kernels,
as in the explicit $n=4$ diagram in E.30.
The resulting mapping-class action preserves the
permutation of every generator, because $\lambda$ is
odd and the conjugating evaluations of $f$ are pure.
On $\widehat B_n^{\rm ab}=\widehat\mathbb Z$ it acts
as multiplication by $\lambda$.

We use the precise finite-subgroup conclusion of
Boggi \cite[Theorem C(i), equivalently Theorem 5.1(i)
with empty boundary]{ERootBoggi2026}.  For the genus-zero
surface with $n+1$ punctures, the natural inclusion of
the discrete full mapping class group in its profinite
completion induces a bijection on conjugacy classes
of solvable finite subgroups.  The surface is connected,
oriented, hyperbolic and of finite type, and the subgroups
used here are finite cyclic groups.  Thus all the hypotheses
are satisfied.  Applying this assertion to the images
of $\langle\delta_n\rangle$ and
$\langle\varepsilon_n\rangle$ in the central quotient
shows that their images under $F$ are conjugate to
discrete cyclic subgroups.

The discrete rotation classification
\cite[Theorem 3.1 and Corollary 3.2]{ERootCollas2012}
then identifies the possible generators.  In the
$(n+1)$-punctured sphere, $\delta_n$ has order $n$
and exactly one marked fixed pole, namely $n+1$.
The element $\varepsilon_n$ has order $n-1$ and
exactly two marked fixed poles.  The other marked
points lie in the rotation orbits.  Powers of these
rotations with the same order have the same indicated
fixed-point pattern, so the preserved permutation
selects the corresponding family.  We obtain
conjugacies to $\delta_n^a$ and $\varepsilon_n^b$,
where $a$ is a unit modulo $n$ and $b$ is a unit
modulo $n-1$.

For $\delta_n$, the unique marked fixed pole forces
the conjugating permutation to fix $n+1$.  The
conjugator therefore belongs to the stabilizer
in \textup{(E.39.5)}.  For $\varepsilon_n$, the
conjugator preserves the set of the two fixed poles.
In the rotation model whose marked points are
$\{0,\infty\}\cup\mu_{n-1}$, the actual holomorphic
map $w(t)=1/t$ exchanges the poles and conjugates
the rotation to its inverse.  Compose with $w$
when the poles are exchanged.  The resulting
conjugator fixes $n+1$, with exponent either $b$
or $-b$; denote that exponent by $b'$.
This operation supplies the required conjugator
inside the stabilizer before determining its exponent.

The abelianization of the stabilizer is
$\mathbb Z/n(n-1)$ by \textup{(E.39.4)}.
The classes of $\delta_n$ and $\varepsilon_n$
are $n-1$ and $n$, respectively.  The actual
action on this quotient therefore gives
\[
 (n-1)a=(n-1)\lambda\pmod{n(n-1)},\qquad
 nb'=n\lambda\pmod{n(n-1)}.
 \tag{E.39.6}
\]
Hence $a\equiv\lambda\pmod n$ and
$b'\equiv\lambda\pmod{n-1}$.  Both central-quotient
conjugacies have the required exponent $\lambda$.
Here $\varepsilon_n^{n-1}=z_n$ in $B_n$:
its rotation has order $n-1$ in \textup{(E.39.4)},
so that power is a boundary-twist power; its
exponent sum $n(n-1)$ makes the power exactly one.

Lift either stabilizer conjugator to $\widehat B_n$.
For $\beta=\delta_n$ or $\varepsilon_n$ the equality
now has the exact form
\[
 F(\beta)=h\beta^\lambda h^{-1}z_n^k,
                         \qquad k\in\widehat\mathbb Z.
 \tag{E.39.7}
\]
Apply the full exponent-sum map to
$\widehat\mathbb Z$.  Both $F(\beta)$ and
$\beta^\lambda$ have exponent sum
$\lambda\operatorname{exp}(\beta)$, whereas
$z_n^k$ has exponent sum $n(n-1)k$.
Multiplication by the nonzero integer $n(n-1)$
is injective on
$\widehat\mathbb Z=\prod_p\mathbb Z_p$.
It follows that $k=0$, proving \textup{(E.39.2)}.
For a discrete conjugate $u\beta^j u^{-1}$,
apply $F$ and use the actual conjugator
$F(u)hu^{-1}$ to obtain the same assertion.
Finally applying the stated finite map $q$ gives
\textup{(E.39.3)} with conjugators $q(h_\delta)$
and $q(h_\varepsilon)$.  This last step uses
the two homomorphisms $q$ and $qF$ from the
same full source, and requires no induced
endomorphism of the fixed finite quotient.
\end{proof}
